\chapter[Dependent inference and causality]{Dependent observations, selective inference, and causality}
\label{ch:dependent-inference}
Time bins from one animal are typically correlated.
Repeating a test at many channels, lags and frequency bands
also creates selection effects. A valid mathematical result
must state which dependence and selection operations it covers.

\section{Variance of a time-series average}
For a zero-mean weakly stationary time series with
$\gamma(k)=\Cov(X_t,X_{t+k})$, direct expansion gives
\[
\Var(\bar X_T)=\frac{1}{T}
\left[\gamma(0)+2\sum_{k=1}^{T-1}
\left(1-\frac{k}{T}\right)\gamma(k)\right].
\]
This \emph{finite-sample identity} needs no central limit theorem.
When $\sum_k|\gamma(k)|<\infty$, the asymptotic variance is
$[\gamma(0)+2\sum_{k\ge1}\gamma(k)]/T$.
For normalized autocorrelation $\rho(k)$, an approximate effective
sample size is
$T_{\mathrm{eff}}=T/[1+2\sum_{k\ge1}\rho(k)]$,
provided that series is meaningful. Negative correlations can make
this heuristic exceed $T$; it must not be blindly clipped or treated
as an exact count of independent animals.

\section{HAC and resampling under dependence}
A heteroskedasticity-and-autocorrelation-consistent (HAC)
variance estimator uses weighted autocovariances, for example
\[
\widehat\Omega=\widehat\gamma(0)
 +2\sum_{k=1}^{L}\left(1-\frac{k}{L+1}\right)
 \widehat\gamma(k).
\]
This is a Bartlett-window Newey--West form in a scalar,
mean-estimation setting. Consistency requires hypotheses on
stationarity, weak dependence, moments and bandwidth growth;
finite samples and nonstationary trials require scrutiny.
Block bootstrap preserves within-block temporal structure
but is only justified under suitable dependence assumptions
and a suitable block length. Randomly permuting every LFP
sample is generally \emph{not} an exchangeable null.

\section{Multiple comparisons and post-selection}
Holm's step-down procedure controls family-wise error under
valid marginal $p$-values without an independence assumption.
Benjamini--Yekutieli controls FDR under arbitrary dependence
using the harmonic correction $c_m=\sum_{j=1}^m1/j$
in its standard formulation; it may be conservative.
A maximum-statistic permutation test controls a family of tests
if the randomization acts on the units exchangeable under the
joint null. Cluster-permutation procedures use adjacency
and a cluster-level statistic; rejection does not locate the
exact significant onset, offset, or single electrode.
After selecting the largest spectral peak and testing it
on the same data, ordinary unadjusted $p$-values generally
fail to reflect that selection. Sample splitting, independent
replication or mathematically justified selective inference
are possible remedies.

\section{Observational dependence versus intervention}
A directed acyclic graph (DAG) encodes a collection of
conditional-independence relations under the Markov property.
Paths are blocked or unblocked according to \emph{d-separation}.
If $Z$ is an appropriate backdoor adjustment set for a proposed
causal effect $X\to Y$, under causal sufficiency assumptions
relative to the DAG and positivity,
\[
p(y\mid\operatorname{do}(x))
=\int p(y\mid x,z)p(z)\,dz.
\]
The formula fails for invalid adjustment sets, unmeasured
confounding, or violated structural assumptions. In particular,
a lagged correlation does not supply the DAG.
Faithfulness is an additional assumption (and can fail by
parameter cancellation); the common-cause principle motivates
but does not uniquely identify latent confounders.

\section{Transfer entropy and Gaussian Granger}
Transfer entropy from $X$ to $Y$, conditional on selected
histories $Y^-$ and $Z^-$, is a conditional mutual information:
$TE=I(Y_t;X^-\mid Y^-,Z^-)$.
For jointly Gaussian variables, with unrestricted correct
linear conditional predictors and covariance determinants
specified consistently,
\[
TE_{X\to Y\mid Z}=\tfrac12
\log\frac{\det\Sigma_{\mathrm{reduced}}}
                {\det\Sigma_{\mathrm{full}}}
\quad\text{(nats)}.
\]
A common multivariate Granger measure is the full log ratio;
therefore it equals $2TE$ in this Gaussian setting.
This is a theorem about \emph{predictive distributions}, not an
identification theorem for causal interventions.
Spectral Granger additionally requires consistent VAR transfer
functions or spectral factorization (e.g., Wilson methods) and
model residual checks. The present
\path{granger_causality.m} computes a temporal
log-variance ratio, not spectral Granger; the existing
\path{transfer_entropy.m} is a first-order,
discrete plug-in conditional-information estimator.
\begin{caution}The two MATLAB primitives are exploratory and do not by themselves establish causal identifiability.\end{caution}